\begin{proposition}[B.3.8; Proved; analytic unramified torsors do not evade the Hodge obstruction]\label{push:B38}
Let $V\to C$ be any compact smooth analytic genus-one surface
locally isomorphic, as a torsor over every base disc, to the
Jacobian surface $\mathscr E\to C$ of B.3.4. Algebraicity
or finite order of this analytic torsor class is not assumed.
Then $b_1(V)=6$, and $V$ is K\"ahler. For every holomorphic
line bundle $M$ on $V$, the projective bundle
$\mathbb P_V(\mathcal O\oplus M)$ is K\"ahler. Every smooth
compact threefold bimeromorphic to such a bundle satisfies
the $\partial\bar\partial$-lemma.

Consequently the necessary non-$\partial\bar\partial$
normalization component in B.4.5 cannot be obtained simply
by replacing the projective torsor of B.3.6 by a nonalgebraic
locally trivial analytic torsor, followed by a ruled bundle
and bimeromorphic modifications. This excludes that specified
repair; it does not assert that every possible cusp
normalization has this form.
\end{proposition}
\begin{proof}
Write $p:\mathscr E\to C$ for the projective Jacobian.
The calculation in B.3.4 gives
$R^1p_*\mathcal O_{\mathscr E}=\mathcal O_C(-\infty)$,
and $g(C)=3$. The coherent Leray sequence therefore gives
$h^1(\mathscr E,\mathcal O_{\mathscr E})=3$. Projectivity
and Hodge decomposition imply $b_1(\mathscr E)=6$.
The topological Leray sequence with rational coefficients is
\[
 0\longrightarrow H^1(C;\mathbb Q)
 \longrightarrow H^1(\mathscr E;\mathbb Q)
 \longrightarrow H^0(C,R^1p_*\mathbb Q)
 \longrightarrow H^2(C;\mathbb Q)
 \longrightarrow H^2(\mathscr E;\mathbb Q).
\]
The last arrow is injective because $p$ has a section.
The first arrow is an isomorphism of six-dimensional
vector spaces. Exactness thus implies
$H^0(C,R^1p_*\mathbb Q)=0$.

The local torsor identifications differ by translations.
On a smooth elliptic fibre translations act trivially on
$H^1$. On the $I_{12}$ fibre they preserve the orientation
of the component polygon and therefore act trivially on
its one-dimensional $H^1$ as well. These identifications
commute with the local specialization maps. Hence, for
$\pi:V\to C$, they glue to
$R^1\pi_*\mathbb Q\simeq R^1p_*\mathbb Q$ on the whole
curve, including infinity. The same Leray edge sequence
for $V$ now gives
$H^1(V;\mathbb Q)\simeq H^1(C;\mathbb Q)$.
Thus $b_1(V)=6$ is even. Buchdahl's theorem
\cite[Theorem 11, p.297]{Buchdahl1999} proves that $V$
is K\"ahler without an algebraicity hypothesis.

Choose a Hermitian metric on $\mathcal O\oplus M$.
The induced metric on the tautological relatively ample
line of its projectivization has a closed real $(1,1)$
curvature form $\theta$ positive in the vertical direction.
If $\omega$ is a K\"ahler form on $V$, compactness implies
that $\theta+A\pi^*\omega$ is positive in every direction
for a sufficiently large constant $A$. It is closed, so
it is a K\"ahler form on the projective bundle. A smooth
compact bimeromorphic model is therefore in Fujiki class
$\mathcal C$ and satisfies the $\partial\bar\partial$-lemma;
the domination statement is recalled, with its precise
reference to the $\partial\bar\partial$ theorem, in
\cite[Section 2.1, p.8]{Chen2026HodgeRoot}. Apply B.4.5
for the final obstruction.
\end{proof}
